\chapter{Ngủ Bù Tiết Đầu}
\chapterintrobi{Tối thức quá khuya nên sáng gửi nợ sang lớp học}{Sleep through the first period}{Đây là thói quen biến tiết đầu thành phần nối dài của chiếc gối ở nhà.}

\begin{lucbatbox}{Lục bát: Đầu giờ còn đang đêm}
Sáng nay em tới rất ngoan\
Chỉ là tinh thần chưa kịp theo sang\
Mắt kia khép mở dịu dàng\
Như còn thương nhớ chiếc chăn tối rồi\
Tiết đầu thành chỗ nghỉ ngơi\
Bài giảng đi trước, em ngồi phía sau
\end{lucbatbox}

\begin{tutuyetbox}{Thất ngôn tứ tuyệt: Nợ giấc ngủ lên trường}
Đêm qua em thức với màn hình\
Sáng nay vào lớp mắt lung linh\
Nếu ngủ đủ trước khi đi học tới\
Tiết đầu đâu phải hóa phòng riêng
\end{tutuyetbox}

\begin{ghichu}
Ngủ bù tiết đầu thường bắt đầu từ thói quen thức khuya không thương ngày mai của chính mình.
\end{ghichu}

\begin{englishnote}
\textbf{first period}: tiết đầu.\par
\textbf{catch up on sleep}: ngủ bù.\par
\textbf{I slept too late last night.}: Tối qua em ngủ quá muộn.
\end{englishnote}